%% chapters/03-convergence.tex
\chapter{Convergence: How Sequences of Random Things Settle Down}
\label{ch:convergence}

\whereWeAre{%
A single random variable is one thing; a \emph{sequence} of random variables is the entry point into asymptotic statistics. Three different notions of ``settling down'' coexist, almost surely, in probability, in distribution, with a clean implication chain between them.%
}

\PLACEHOLDER{Sources: existing main.tex `Modes of Convergence' section; Joseph thesis Section 2.2; Casella \& Berger Section 5.5.}

\section{Almost sure convergence}
\PLACEHOLDER{Definition, sample-path interpretation, examples (shifted RVs, $X + 1/n$).}

\section{Convergence in probability}
\PLACEHOLDER{Definition, contrast with a.s.\, the canonical Borel--Cantelli counterexample.}

\section{Markov's and Chebyshev's inequalities}
\PLACEHOLDER{Statements, proofs, used in the WLLN.}

\section{The Weak Law of Large Numbers}
\PLACEHOLDER{Statement, Chebyshev-based proof.}

\section{Convergence in distribution}
\PLACEHOLDER{Definition via CDFs, contrast with the other modes, why it's the right notion for the CLT.}

\section{The implication chain}
\PLACEHOLDER{a.s.\ $\Rightarrow$ in prob $\Rightarrow$ in dist; converses fail.}

\section{Brief warm-ups}
\PLACEHOLDER{3--5 short questions with immediate answers.}

\chapterRecap{%
\begin{itemize}
  \item \PLACEHOLDER{Three modes, ordered from strongest to weakest.}
  \item \PLACEHOLDER{Markov + Chebyshev are the technical workhorses.}
  \item \PLACEHOLDER{Convergence in distribution is what limit theorems are about.}
\end{itemize}%
}

\section*{Exercises}
\PLACEHOLDER{8--15 longer exercises.}
